% !TEX program = pdflatex
% SafeBoxForest v5  T-RO Journal Extension
% Compile: pdflatex -> bibtex -> pdflatex x 2
\documentclass[journal]{IEEEtran}

% --- Packages ---
\usepackage[utf8]{inputenc}
\usepackage[T1]{fontenc}
\usepackage{amsmath,amssymb,amsthm}
\usepackage{graphicx}
\usepackage{algorithm}
\usepackage[noend]{algpseudocode}
\def\Or{\textbf{or}}
\newcommand{\alAnd}{\textbf{and}}
\usepackage{booktabs}
\usepackage{multirow}
\usepackage{cite}
\usepackage{xcolor}
\usepackage{hyperref}
\usepackage{cleveref}
\usepackage{bm}
\usepackage{tabularx}
\usepackage{xspace}
\usepackage{microtype}

% --- Theorem environments ---
\newtheorem{theorem}{Theorem}
\newtheorem{definition}{Definition}
\newtheorem{remark}{Remark}
\newtheorem{proposition}{Proposition}
\newtheorem{assumption}{Assumption}

% --- Custom commands ---
\newcommand{\Cfree}{\mathcal{C}_\text{free}}
\newcommand{\Cspace}{\mathcal{C}}
\newcommand{\AABB}{\text{AABB}}
\newcommand{\iAABB}{\text{iAABB}}
\newcommand{\FK}{\text{FK}}
\newcommand{\SBF}{\text{SBF}}
\newcommand{\Real}{\mathbb{R}}
\newcommand{\calO}{\mathcal{O}}
\newcommand{\calA}{\mathcal{A}}
\newcommand{\calF}{\mathcal{F}}
\newcommand{\calR}{\mathcal{R}}
\newcommand{\bT}{\mathbf{T}}
\newcommand{\TODO}[1]{\textit{[TODO: #1]}}
\newcommand{\ie}{\textit{i.e.},\xspace}
\newcommand{\eg}{\textit{e.g.},\xspace}
\newcommand{\cf}{\textit{cf.}\@\xspace}
\newcommand{\etal}{\textit{et~al.}\@\xspace}
\newcommand{\LinkIAABBDef}{Per-link axis-aligned box obtained from the bounding box of a link's endpoint-iAABB pair, inflated by link radius; when subdivision is enabled, it denotes the union of the subdivided sub-link boxes.}
\DeclareMathOperator{\interior}{int}

% ==============================================================================
\begin{document}

\title{SafeBoxForest: Rapid C-space Box Cover with\\
Lifelong Envelope Caching for GCS Planning}

\author{
  TIAN,~Yu and Ren,~Hongliang
  \thanks{Department of Electronic Engineering,
  The Chinese University of Hong Kong, Hong Kong SAR, China.
  \{ty1997, hlren\}@ee.cuhk.edu.hk}
}

\markboth{IEEE Transactions on Robotics, Vol.~X, No.~Y, 2026}{}

\maketitle

% ==============================================================================
% ABSTRACT
% ==============================================================================
\begin{abstract}
We present SafeBoxForest~(SBF), a system that rapidly
covers the collision-free configuration space ($\Cfree$) of
articulated robots with convex boxes
for Graph-of-Convex-Sets~(GCS) motion planning.
SBF introduces a modular two-layer pipeline:
four interchangeable \emph{endpoint-iAABB sources}
(IFK, CritSample, AnalyticalCrit, GCPC)
feed into three \emph{envelope representations}
(LinkIAABB\footnotemark, LinkIAABB\_Grid, Hull16\_Grid).
All sources produce a unified \emph{endpoint-iAABB} intermediate
representation---per-endpoint interval axis-aligned bounding boxes
along the kinematic chain---which decouples
source computation from envelope rasterisation.
Key components include:
(i)~a \emph{five-phase analytical critical-point solver}
with DH-chain direct coefficient extraction (zero-FK optimisation)
and interval-FK certification
(completeness for face-$k$, $k \leq 2$);
(ii)~a \emph{Geometric Critical-Point Cache}~(GCPC) that
pre-tabulates FK-evaluated critical configurations via an
eight-phase query pipeline with $k\pi/2$-background caching,
incremental DH-chain recomputation, and AA pruning,
answering per-box queries $1.1{\times}$ faster than the
analytical solver;
(iii)~a lightweight \emph{CritSample} that enumerates
boundary $k\pi/2$ critical angles with precomputed DH
transforms for fast critical-point sampling
(no GCPC cache required);
(iv)~the \emph{Sparse Voxel BitMask}---an $8^3$ BitBrick
tile map with turbo scanline rasterisation---producing
hull-16 envelopes $35\%$ tighter than axis-aligned bounding
boxes with constant-time collision queries; and
(v)~the \emph{Lifelong Envelope Cache Tree}~(LECT)
with scene pre-rasterisation ($2272{\times}$ collision speedup)
and zero-loss voxel merging.
We benchmark all $4{\times}3{=}12$ pipeline configurations on a
7-DOF manipulator.
The recommended configuration (CritSample+Hull16\_Grid)
achieves $0.117$\,m$^3$ envelope volume at $0.278$\,ms
cold-start cost.
The AnalyticalCrit and GCPC sources validate CritSample
accuracy: their Hull16\_Grid volumes ($0.121$--$0.123$\,m$^3$) differ
by only~$1.4\%$, confirming that boundary sampling captures
virtually all extrema.
Among the four sources, only IFK produces \emph{certified}
(provably collision-free) envelopes via conservative
interval over-approximation;
Analytical and GCPC compute tighter envelopes that are
safe in the vast majority of cases but do not carry
formal guarantees, while CritSample trades a small
additional completeness gap for sub-millisecond latency.
\end{abstract}
\footnotetext{\LinkIAABBDef}

\begin{IEEEkeywords}
Motion planning, GCS, convex decomposition, interval arithmetic,
collision-free regions
\end{IEEEkeywords}

% ==============================================================================
\section{Introduction}\label{sec:intro}
% ==============================================================================

Motion planning for manipulators in cluttered environments
remains a fundamental challenge.
Sampling-based planners (RRT, PRM) offer general-purpose
solutions but provide no certificates and scale poorly in
narrow passages~\cite{lavalle2006planning}.
Recent work on GCS-based planning~\cite{marcucci2024motion}
formulates trajectories as convex programs over
safe regions, guaranteeing optimality and smoothness---provided
the regions accurately cover $\Cfree$.

The bottleneck shifts to \emph{region generation}.
IRIS-NP~\cite{werner2024fast} inflates ellipsoidal regions
via semidefinite programming but requires $0.3$--$1$\,s per
region and does not reuse geometry across queries.
Clique covers~\cite{amice2024clique} offer better coverage but
are limited to polyhedral obstacle representations.

We propose SafeBoxForest~(SBF), which decomposes $\Cfree$
into axis-aligned convex boxes.
When using IFK as the endpoint source, each box's collision-free
status is \emph{certified} through conservative interval
over-approximation; the tighter Analytical and GCPC sources
produce envelopes that are empirically safe but lack
formal guarantees.
Key to the approach is a hierarchy of envelope
representations---from cheap axis-aligned bounding boxes to
tight sparse-voxel hull-16 grids---stored in a persistent
Lifelong Envelope Cache Tree~(LECT) that amortises computation
across boxes.

\textbf{Contributions.}
\begin{itemize}
  \item A five-phase analytical critical-point solver with
    DH-chain direct coefficient extraction (zero-FK
    optimisation) that computes provably tight envelopes with
    face-$k$ completeness for $k \leq 2$
    (\cref{sec:critical}).
  \item An eight-phase Geometric Critical-Point Cache (GCPC)
    with $q_0$ elimination, $q_1$ period-$\pi$ symmetry,
    $k\pi/2$-background caching, incremental chain
    recomputation, and multi-stage AA pruning for
    $1.1{\times}$ faster per-box queries than the analytical solver,
    validating CritSample accuracy (\cref{sec:gcpc}).
  \item A lightweight CritSample design using boundary
    $k\pi/2$ enumeration with precomputed DH transforms,
    delivering near-analytical quality at sub-millisecond
    cost without requiring the GCPC cache
    (\cref{sec:critsample}).
  \item The Sparse Voxel BitMask representation with turbo
    scanline rasterisation and constant-time collision/merge
    operations (\cref{sec:voxel}).
  \item The LECT data structure with scene pre-rasterisation
    ($2272{\times}$ collision speedup) and hull-based
    zero-loss coarsening (\cref{sec:sbf}).
  \item A comprehensive $4{\times}3$ pipeline benchmark
    demonstrating the effectiveness of each component on
    a 7-DOF manipulator (\cref{sec:experiments}).
\end{itemize}

% ==============================================================================
\section{Related Work}\label{sec:related}
% ==============================================================================

\paragraph{Convex decomposition for motion planning.}
VCD~\cite{werner2024vcd} decomposes task space into convex polytopes
via vertex enumeration, but the combinatorial complexity grows
exponentially in dimension.
IRIS~\cite{deits2015computing} and its recent extension
IRIS-NP~\cite{werner2024fast} iteratively inflate ellipsoids
in $\Cfree$ using SDP or nonlinear programming;
each region costs $0.3$--$1$\,s and is not reusable across scenes.
Clique covers~\cite{amice2024clique} partition $\Cfree$ using
graph-theoretic methods, achieving good coverage but requiring
explicit facet representations.

\paragraph{GCS motion planning.}
Marcucci et al.~\cite{marcucci2024motion} formulates motion
planning as a shortest-path problem over a graph where
vertices are convex safe regions and edges encode transitions.
The resulting convex programs guarantee global optimality
(given the regions) and produce smooth trajectories.
SBF provides the convex regions required by GCS.

\paragraph{Interval arithmetic in robotics.}
Interval analysis has been used for verified collision
checking~\cite{merlet2004solving}, workspace
computation~\cite{merlet2006interval}, and motion
planning~\cite{jaulin2001applied}.
Our work exploits interval FK to produce \emph{certified}
bounding volumes without point-sampling, extending the
approach to multi-layer envelope representations.

\paragraph{Spatial hashing and voxel methods.}
Spatial hashing dates to Teschner et al.~\cite{teschner2003optimized}
for deformable body collision.
Recent work uses hierarchical voxel grids (VDB~\cite{museth2013vdb})
for robotic occupancy mapping.
Our BitBrick layout is lighter-weight: a single-level
$8^3$ tile map that trades hierarchy for cache-friendly
bitwise collision and zero-loss Boolean merge.

% ==============================================================================
\section{Certified Envelope Computation}\label{sec:envelope}
% ==============================================================================

% Code: include/sbf/robot/interval_fk.h, include/sbf/envelope/endpoint_source.h,
%       include/sbf/envelope/analytical_solve.h, include/sbf/envelope/analytical_coeff.h

This section presents the mathematical foundations for computing
provably conservative workspace envelopes from C-space boxes.

\subsection{Interval Forward Kinematics}\label{sec:ifk}
% Code: compute_fk_full() in robot/interval_fk.h
% Code: I_sin(), I_cos(), build_dh_joint(), imat_mul_dh() in robot/interval_math.h
% Code: extract_link_iaabbs() in robot/ifk_utils.h

Consider a $D$-DOF serial manipulator with $L$ collision-relevant links.
Given a C-space box $B \subset \Real^D$ and obstacle iAABBs $\calO$,
the goal is to certify $B \subseteq \Cfree$.

\emph{Link-iAABB envelope.}
The \emph{link-iAABB envelope} $\iAABB_\ell(B)$ satisfies
\begin{equation}\label{eq:envelope}
  \iAABB_\ell(B) \supseteq \bigcup_{\bm{q} \in B}
    \mathrm{geom}_\ell\bigl(\FK(\bm{q})\bigr).
\end{equation}
Each link~$\ell$ is modeled as a capsule: a line segment between joint
origins $\bm{p}_\ell^{\mathrm{prox}}$ and $\bm{p}_\ell^{\mathrm{dist}}$,
inflated by radius~$r_\ell$.

\begin{assumption}[Geometry Enclosure]\label{asm:geom}
For each link~$\ell$, the capsule of radius~$r_\ell$ encloses the
true mesh geometry at every configuration.
\end{assumption}

Since every interior point of the capsule segment is a convex combination
of its two endpoints, coordinate extremes over all configurations in~$B$
are attained at the endpoints alone.
The iAABB therefore reduces to
\begin{equation}\label{eq:endpoint_iaabb}
  \iAABB_\ell(B)
  = \mathrm{bbox}\!\Bigl(
      \bigl[\bm{p}_\ell^{\mathrm{prox}}\bigr]_B,\;
      \bigl[\bm{p}_\ell^{\mathrm{dist}}\bigr]_B
    \Bigr) \oplus r_\ell,
\end{equation}
where $[\cdot]_B$ is the interval enclosure over $B$ and $\oplus\,r_\ell$
is a uniform inflation.

\emph{Computing via interval-FK.}
Replacing scalar joint angles with interval counterparts in the DH chain
and multiplying sequentially,
\begin{equation}\label{eq:chain}
  [\bT^{0:k}] = [\bT^{0:k-1}] \otimes [\bT_k],
\end{equation}
yields interval transforms for all endpoints.
The two endpoint positions in~\eqref{eq:endpoint_iaabb} are read directly
from the translation components of $[\bT^{0:\ell-1}]$ and $[\bT^{0:\ell}]$;
no additional FK evaluations are needed.
By inclusion-monotonicity of interval
arithmetic~\cite{moore2009introduction}, the result is a
provably conservative enclosure; chain multiplication may over-approximate
(the \emph{wrapping effect}), affecting tightness but not soundness.
By construction, over-approximation can only cause \emph{false rejections}
(a free box mistakenly rejected), never false acceptances, so every
box that passes the envelope check is guaranteed collision-free.

\emph{Active-link filtering.}
Not all links contribute to collision geometry.
For each link~$\ell$, we define its \emph{inherent length}
$L_\ell = \sqrt{a_\ell^2 + d_\ell^2}$ from the DH
parameters~$(a_\ell, d_\ell)$, which is independent of the joint
angle.
Link~0 (the base frame) is always excluded.
Links with $L_\ell < \tau$ (default $\tau = 0.1\,\text{m}$) are
treated as zero-length and excluded, as their near-degenerate
iAABBs add no useful collision information.
Only $L_\text{act} \leq L$ \emph{active links} are retained, reducing
both storage and collision-check cost.
The active-link map and associated metadata (per-link collision
radii, affecting-joint counts, and an iFK cutoff frame) are
determined once from the robot model and stored
in a self-contained \texttt{IAABBConfig} structure,
which can also be serialised to JSON for offline reuse.
% Code: IAABBConfig::from_robot() in robot/iaabb_config.h

\subsection{Subdivision Enhancements}\label{sec:subdivision}
% Code: derive_iaabb_subdivided() in envelope/envelope_derive.h
% Code: compute_fk_incremental() for prefix reuse

Interval-FK can over-approximate due to the wrapping effect.
Two complementary subdivision strategies reduce this conservatism while
preserving certification.

\emph{C-space box subdivision.}
Bisecting box~$B$ along a dimension into $B_1$ and $B_2$ and computing
their envelopes separately yields tighter enclosures than treating $B$
as a whole.
Recursive $k$-level subdivision produces $2^k$ sub-boxes.
Certification is preserved by inclusion-monotonicity:
each sub-box envelope is a valid over-approximation, and their union
is contained in the original ($\bigcup_i \iAABB_\ell(B_i) \subseteq \iAABB_\ell(B)$)
but pointwise tighter.
Sub-boxes share FK prefix chains up to the split dimension,
so the extra cost is sub-linear.

\emph{Incremental FK.}
When a parent node splits along dimension~$d$, the interval transform
prefix $[\bT^{0:d-1}]$ is identical for both children.
Only the suffix $[\bT^{d:D}]$ must be recomputed.
On a 7-DOF manipulator, this saves approximately 43\% of FK cost on
average; the saving grows with joint index~$d$.

\emph{Link subdivision.}
For a long link whose iAABB spans a large workspace volume,
splitting link~$\ell$ into $n$ equal segments and computing per-segment
iAABBs independently yields a tighter aggregate envelope.
The $n{-}1$ intermediate segment endpoints are obtained by linearly
interpolating the proximal and distal joint-origin positions
$\bm{p}_\ell^{\mathrm{prox}}$ and $\bm{p}_\ell^{\mathrm{dist}}$,
which are already available from the FK chain; no additional FK
evaluation is required.
Certification is preserved: each sub-segment iAABB encloses the
corresponding mesh portion under \cref{asm:geom}.
Accordingly, the certified occupied set of a subdivided link is the
\emph{union} of these sub-segment iAABBs, not the full interior of the
coarse whole-link AABB. For slender or diagonal links, interior points of
the coarse AABB may lie in empty space and need not be covered by any
individual sub-segment box.

\subsection{Endpoint-iAABB: Unified Intermediate Representation}%
\label{sec:endpoint_iaabb}

All envelope derivation methods produce a common intermediate
representation called the \emph{endpoint-iAABB}.
For each endpoint $e \in \{0,1,\ldots,n_\text{endpoints}-1\}$
(where $n_\text{endpoints} = n_\text{joints} + \mathbf{1}_\text{tool}$),
the endpoint-iAABB is the interval axis-aligned bounding box
\begin{equation}\label{eq:ep_iaabb}
  E_{e} \;=\;
  \bigl[\mathrm{lo}_x,\,\mathrm{lo}_y,\,\mathrm{lo}_z,\;
        \mathrm{hi}_x,\,\mathrm{hi}_y,\,\mathrm{hi}_z\bigr],
\end{equation}
where each $[\mathrm{lo}_d,\mathrm{hi}_d]$ is the interval enclosure
of the $d$-th coordinate of the endpoint over the C-space
box~$B$.
The total storage is $n_\text{endpoints} \times 6$ floats.

The key property is that the envelope representation (LinkIAABB,
LinkIAABB\_Grid, or Hull16\_Grid) is derived \emph{solely} from
the endpoint-iAABB array---it does not depend on how the
endpoint-iAABBs were produced.
This decouples the \emph{source} computation
(IFK, CritSample, AnalyticalCrit, or GCPC) from the
\emph{envelope rasterisation}:
\begin{equation}\label{eq:pipeline}
  \begin{aligned}
    &\underbrace{\text{Source}}_{\text{IFK / CritSample / Analytical / GCPC}}
    \;\xrightarrow{\;E_e\;} \\
    &\underbrace{\text{Envelope}}_{\text{LinkIAABB / LinkIAABB\_Grid / Hull16\_Grid}}.
  \end{aligned}
\end{equation}

\emph{Deriving link-iAABBs from endpoint-iAABBs.}
For each active link~$\ell$ with parent endpoint $e_{\ell-1}$
and child endpoint $e_\ell$, the link-iAABB is
\begin{equation}
  \iAABB_\ell = \mathrm{bbox}\bigl(E_{e_{\ell-1}},\, E_{e_\ell}\bigr)
  \;\oplus\; r_\ell,
\end{equation}
where $\oplus\,r_\ell$ is per-link radius inflation.
This is computed by \texttt{extract\_link\_iaabbs()}.

\emph{Deriving Hull16\_Grid from endpoint-iAABBs.}
The proximal box $E_{e_{\ell-1}}$ and distal box $E_{e_\ell}$
define the 16-vertex convex hull of \cref{sec:hull16}.
The hull-16 turbo rasteriser takes these two endpoint iAABBs
(plus link radius) directly, without passing through the
link-iAABB intermediate.

\subsection{Analytical Critical-Point Solver}\label{sec:critical}
% Code: include/sbf/envelope/analytical_solve.h
% Code: include/sbf/envelope/analytical_coeff.h
% Code: include/sbf/envelope/analytical_utils.h

Interval-FK provides a \emph{certified} envelope but may over-approximate
due to the wrapping effect.
We complement it with a \emph{tight} envelope obtained by solving for
the true coordinate extremes analytically.
We first establish structural properties of DH kinematics that
enable exact solutions, then present the 5-phase solver and prove
its completeness guarantees.

% ---------- structural properties ----------

\subsubsection{Trigonometric Structure of DH Chains}\label{sec:critical_structure}

Each $4{\times}4$ homogeneous DH matrix $\bT_j$ for a revolute joint~$j$
is affine in $(\cos q_j,\,\sin q_j)$:
\begin{equation}\label{eq:dh_affine}
  \bT_j(q_j)
  = \bm{M}_j^{(c)}\cos q_j
  + \bm{M}_j^{(s)}\sin q_j
  + \bm{M}_j^{(0)},
\end{equation}
where $\bm{M}_j^{(c)}, \bm{M}_j^{(s)}, \bm{M}_j^{(0)}$ are constant
matrices determined by the DH parameters $(a_j,\alpha_j,d_j)$.
The composed transform is
$\bT^{0:\ell} = \bT_1 \cdots \bT_\ell$
and the position coordinate of the distal origin is
$p_{\ell,d}(\bm{q}) = [\bT^{0:\ell}]_{d,4}$.
Since each factor is affine in its own $(c_j,s_j)$ pair,
the product is \emph{multilinear} in
$\{(\cos q_j,\sin q_j)\}_{j=1}^\ell$.

\begin{proposition}[Univariate sinusoidal]\label{prop:sinusoid}
Fix all joints except~$j$.
Then
\begin{equation}\label{eq:sin_form}
  p_{\ell,d}(q_j)
  = A_j\cos q_j + B_j\sin q_j + C_j,
\end{equation}
where $A_j, B_j, C_j$ depend only on the fixed joints and
DH parameters.
This function has \emph{exactly two} critical points per $2\pi$
period: $q_j^* = \mathrm{atan2}(B_j,A_j)$ (global max)
and $q_j^* + \pi$ (global min).
\end{proposition}
\begin{proof}
By~\eqref{eq:dh_affine}, $\bT_j$ is the only factor containing
$q_j$; all other factors are constant matrices.
Collecting terms in the matrix product yields~\eqref{eq:sin_form}.
The derivative
$\partial p/\partial q_j = -A_j\sin q_j + B_j\cos q_j = 0$
gives $\tan q_j = B_j/A_j$, solved uniquely (mod $\pi$) by $\mathrm{atan2}$.
\end{proof}

\begin{proposition}[Bivariate bilinear trigonometric]\label{prop:bilinear}
Fix all joints except $(i,j)$.
Then
\begin{equation}\label{eq:bilinear}
  p_{\ell,d}(q_i,q_j) = \sum_{m=1}^{9} a_m\,\phi_m(q_i,q_j),
\end{equation}
where $\{\phi_m\} = \{c_i c_j,\; c_i s_j,\; s_i c_j,\; s_i s_j,\;
c_i,\; s_i,\; c_j,\; s_j,\; 1\}$ and the nine coefficients $a_m$
depend on the fixed joints.
\end{proposition}
\begin{proof}
$\bT_i$ and $\bT_j$ are distinct factors in the matrix product
$\bT^{0:\ell}$.
Each contributes $(c_k, s_k)$ linearly; their product introduces
exactly the four cross-terms $c_i c_j, c_i s_j, s_i c_j, s_i s_j$
plus the six lower-order and constant terms.
\end{proof}

\begin{proposition}[Bivariate critical-point resultant]\label{prop:resultant}
Setting $\nabla_{(q_i,q_j)} p_{\ell,d} = \bm{0}$ in the bilinear
model~\eqref{eq:bilinear} and applying the Weierstrass substitution
$t_j = \tan(q_j/2)$ reduces the system to a single univariate
polynomial of degree~$\leq 8$ in~$t_j$.
All real roots can therefore be found exactly via the companion-matrix
eigenvalue method.
\end{proposition}
\begin{proof}
Setting $\partial p / \partial q_i = 0$ yields
$\tan q_i = P(q_j) / Q(q_j)$ where $P,Q$ are of the
form $\alpha c_j + \beta s_j + \gamma$.
Substituting into $\partial p / \partial q_j = 0$ and clearing
denominators via $c_j = (1-t_j^2)/(1+t_j^2)$,
$s_j = 2t_j/(1+t_j^2)$ gives a rational equation whose numerator
is a polynomial in $t_j$ of degree~$\leq 8$.
The companion matrix of this polynomial has eigenvalues at all roots,
including complex ones; only real eigenvalues correspond to
physical configurations.
\end{proof}

% ---------- DH direct coefficient extraction ----------

\subsubsection{DH-Chain Direct Coefficient Extraction}%
\label{sec:dh_coeff}

% Code: analytical_coeff.h — extract_1d_coeffs(), extract_2d_coeffs()
% Code: analytical_utils.h — FKWorkspace, compute_prefix()

Traditional approaches to coefficient extraction
(\cref{prop:sinusoid,prop:bilinear}) evaluate FK at a sampling grid
and solve a least-squares system (e.g., $3{\times}3$ QR for 1D,
$9{\times}9$ for 2D).
We replace this with \emph{direct algebraic extraction} from the
DH chain, eliminating all FK evaluations.

\emph{Prefix--suffix factorisation.}
The composed transform $\bT^{0:\ell}$ factorises as
\begin{equation}\label{eq:prefix_suffix}
  \bT^{0:\ell} = \underbrace{\bT^{0:j-1}}_\text{prefix}
  \;\cdot\; \bT_j(q_j)
  \;\cdot\; \underbrace{\bT^{j+1:\ell}}_\text{suffix}.
\end{equation}
The prefix and suffix are constant (all free joints fixed).
By~\eqref{eq:dh_affine}, only $\bT_j$ depends on~$q_j$,
and the translation column of the product yields the 1D coefficients:
\begin{equation}\label{eq:1d_direct}
  \begin{pmatrix} A_d \\ B_d \\ C_d \end{pmatrix}
  = \text{prefix}_{d,\cdot} \;
  \begin{pmatrix}
    \bm{M}_j^{(c)} \cdot \bm{s}^{(\ell)}_{4} \\
    \bm{M}_j^{(s)} \cdot \bm{s}^{(\ell)}_{4} \\
    \bm{M}_j^{(0)} \cdot \bm{s}^{(\ell)}_{4}
  \end{pmatrix},
\end{equation}
where $\bm{s}^{(\ell)}_4$ is the 4th column of the suffix.
This requires ${\sim}30$ multiplications and zero FK calls.

\emph{2D coefficient extraction.}
For a pair $(i,j)$ with $i < j$, the chain factorises as
$\bT^{0:\ell} = P \cdot \bT_i \cdot M \cdot \bT_j \cdot S$
where $P,M,S$ are constant prefix, middle, and suffix matrices.
The 9 bilinear-trig coefficients of~\eqref{eq:bilinear} are read
from the appropriate columns of $P \cdot [\bm{M}_i^{(c)}, \bm{M}_i^{(s)},
\bm{M}_i^{(0)}] \cdot M \cdot [\bm{M}_j^{(c)}, \bm{M}_j^{(s)}, \bm{M}_j^{(0)}]
\cdot S_{:,4}$,
requiring ${\sim}120$ multiplications and zero FK calls.

\begin{remark}[Speedup]
On a 7-DOF Panda manipulator, replacing QR-based fitting with direct
extraction reduces Phase~1 + Phase~2 coefficient computation from
$22.9$\,ms to $14.4$\,ms ($37.0\%$ speedup).
The prefix matrices are cached in an \texttt{FKWorkspace} structure
across phases, avoiding redundant prefix recomputation.
\end{remark}

% ---------- 5-phase solver ----------

\subsubsection{Five-Phase Solver}\label{sec:five_phase}

We now describe the solver.
Consider maximising (or minimising) $p_{\ell,d}(\bm{q})$ over the box
$B = \prod_{j=1}^{D} [q_j^-,\, q_j^+]$.
By compactness, the extremum exists.
By the Karush--Kuhn--Tucker (KKT) conditions on the box constraints,
at any extremiser~$\bm{q}^*$ the joints partition into
a \emph{free set}
$F = \{j : q_j^* \in (q_j^-,q_j^+)\}$ and a
\emph{boundary set} $\bar F = \{j : q_j^* \in \{q_j^-,q_j^+\}\}$,
with $\partial p_{\ell,d}/\partial q_j(\bm{q}^*) = 0$ for every
$j \in F$.
The five phases enumerate candidate extremisers by face dimension:

\emph{Phase~0 -- Vertex enumeration ($|F|=0$).}
The candidate set for each joint is
$\mathcal{K}_j = \{q_j^-,\, q_j^+\}
\cup \{ k\pi/2 : k\in \mathbb{Z},\;
q_j^- < k\pi/2 < q_j^+ \}$,
and all configurations in the Cartesian product
$\mathcal{K}_1 \times \cdots \times \mathcal{K}_D$ are evaluated.

\emph{Phase~1 -- Edge critical points ($|F|=1$).}
For each joint~$j$ and each of the~$2^{D-1}$ boundary
assignments of the remaining joints,
the 1D coefficients~\eqref{eq:1d_direct} are extracted directly
from the DH chain (via \texttt{extract\_1d\_coeffs}),
and the two critical points
$q_j^* = \mathrm{atan2}(B_j,A_j)$ and $q_j^*+\pi$ are tested.

\emph{Phase~2 -- Face critical points ($|F|=2$).}
For each joint pair $(i,j)$ and each boundary assignment of the
remaining $D{-}2$ joints, the 9 bilinear-trig coefficients
are extracted directly from the DH chain
(via \texttt{extract\_2d\_coeffs}).
The degree-$\leq 8$ resultant (\cref{prop:resultant}) is then
solved via the companion matrix.

\emph{Phase~2.5 -- Pair-constrained critical points.}
\begin{proposition}[Pair-manifold reduction]\label{prop:pair_reduce}
On the affine manifold $q_i + q_j = C$ with $C = k\pi/2$,
substituting $q_j = C - q_i$ into~\eqref{eq:bilinear} and fixing
all other joints reduces $p_{\ell,d}$ to the sinusoidal
form~\eqref{eq:sin_form} in~$q_i$ alone.
\end{proposition}
\begin{proof}
$\cos(C-q_i) = \cos C \cos q_i + \sin C \sin q_i$ and
$\sin(C-q_i) = \sin C \cos q_i - \cos C \sin q_i$ are linear
in $(\cos q_i, \sin q_i)$.
Substituting into every $(c_j, s_j)$ term of~\eqref{eq:bilinear}
collapses the bilinear form to $A'\cos q_i + B'\sin q_i + C'$.
\end{proof}
Phase~2.5 enumerates all integer~$k$ with
$q_i^- + q_j^- \leq k\pi/2 \leq q_i^+ + q_j^+$,
substituting $q_j = k\pi/2 - q_i$,
and solving the resulting sinusoidal exactly as in Phase~1.

\emph{Phase~3 -- Interior coordinate descent ($|F|\geq 3$).}
For faces of dimension~$\geq 3$ the KKT system is a system of
$|F|$ coupled trigonometric equations.
We solve it approximately by block-coordinate descent:
in each sweep, every free joint~$j$ is optimised exactly
(via \cref{prop:sinusoid}) with all others held fixed, cycling
repeatedly.
Multi-start initialisation from Phase~0/1/2/2.5 seeds and
pair-coupled updates improve convergence.

% ---------- completeness ----------

\subsubsection{Completeness Guarantees}\label{sec:completeness}

\begin{theorem}[Face-$k$ completeness for $k \leq 2$]\label{thm:face_complete}
Let $\bm{q}^*$ be a global extremiser of $p_{\ell,d}$ over
$B = \prod_j [q_j^-,q_j^+]$ with free set $F$ satisfying $|F|\leq 2$.
Then the analytical solver finds a configuration with equal or
better objective value, provided that the boundary enumeration
covers all $2^{|\bar F|}$ assignments of $\bar F$.
\end{theorem}
\begin{proof}
$|F|=0$: Phase~0 includes~$\bm{q}^*$.
$|F|=1$: Phase~1 enumerates all boundary assignments and finds
all sinusoidal critical points.
$|F|=2$: Phase~2 enumerates all boundary assignments and finds
all bilinear-trig critical points via the companion matrix.
\end{proof}

\begin{theorem}[Pair-manifold completeness]\label{thm:pair_complete}
Let $\bm{q}^*$ be a global extremiser with
$|F| \leq 2$ and $q_i^* + q_j^* = k\pi/2$ for some integer~$k$.
Then Phase~2.5 finds a configuration with equal or better objective
value.
\end{theorem}

% ---------- fused dual phase 3 ----------

\subsubsection{Fused Dual-Phase-3 Mode}\label{sec:dual_phase3}

Phase~3 is inherently heuristic.
We identify a \emph{basin-shift} problem: Phase~2's companion-matrix
solutions may move the running best~$\bm{q}^\text{best}$ into the basin
of a worse local extremum for coordinate descent.
We solve this by checkpointing per-link extremes after Phase~1
($E^{01}$), proceeding through Phases~2--3 to obtain $E^{0123}$,
running a second Phase~3 from $E^{01}$ to get $E^{013}$,
and merging:
\begin{equation}\label{eq:dual_merge}
  E^*_{\ell,d} = \bigl[\,
    \min\!\bigl(E^{0123}_{\ell,d,\mathrm{lo}},\;
                E^{013}_{\ell,d,\mathrm{lo}}\bigr),\;
    \max\!\bigl(E^{0123}_{\ell,d,\mathrm{hi}},\;
                E^{013}_{\ell,d,\mathrm{hi}}\bigr)
  \,\bigr].
\end{equation}

% ---------- interval-FK certification ----------

\subsubsection{Interval-FK Certification}\label{sec:ifk_cert}

The analytical solver's output satisfies
$\iAABB^{\mathrm{anal}}_\ell \subseteq \iAABB^{\mathrm{true}}_\ell$
(tight but possibly incomplete for $|F| \geq 3$).
Interval-FK produces
$\iAABB^{\mathrm{ifk}}_\ell \supseteq \iAABB^{\mathrm{true}}_\ell$
(conservative outer bound).
The certification step clamps the output:
\begin{equation}\label{eq:cert_clamp}
  \begin{aligned}
    \widehat{\iAABB}_{\ell,d}^{\mathrm{lo}}
      &= \min\!\bigl(\iAABB_{\ell,d}^{\mathrm{anal,lo}},\;
                       \iAABB_{\ell,d}^{\mathrm{ifk,lo}}\bigr), \\
    \widehat{\iAABB}_{\ell,d}^{\mathrm{hi}}
      &= \max\!\bigl(\iAABB_{\ell,d}^{\mathrm{anal,hi}},\;
                       \iAABB_{\ell,d}^{\mathrm{ifk,hi}}\bigr).
  \end{aligned}
\end{equation}

\begin{theorem}[Certified soundness]\label{thm:cert_sound}
With certification enabled, $\widehat{\iAABB}_\ell \supseteq
\iAABB^{\mathrm{true}}_\ell$ for every active link~$\ell$,
regardless of Phase~3 completeness.
\end{theorem}

\subsection{CritSample: Boundary Critical Sampling}%
\label{sec:critsample}
% Code: include/sbf/envelope/crit_sample.h, src/envelope/crit_sample.cpp

CritSample provides near-analytical quality at sub-millisecond cost
via a lightweight boundary-only design that requires
no GCPC cache.

\emph{Critical angle enumeration.}
For each joint~$j$, the critical angle set is
\begin{equation}
  C_j = \{q_j^-,\, q_j^+\}
  \cup \bigl\{ k\tfrac{\pi}{2} : k\tfrac{\pi}{2} \in (q_j^-,q_j^+),\;
  k \in \mathbb{Z} \bigr\}.
\end{equation}
The Cartesian product $C_1 \times \cdots \times C_D$
(bounded by \texttt{max\_boundary\_combos}) is evaluated
using precomputed DH transform matrices.

\emph{Precomputed DH transforms.}
All DH transforms $\bT_j(c) = \bT_{j-1} \cdot \mathbf{DH}_j(c)$ for
each $(j, c \in C_j)$ pair are precomputed once before the
enumeration loop, eliminating all $\sin/\cos$ calls from the
hot path (typically $\sim$500 matrix multiplies per query).
The Cartesian product is traversed via an explicit stack-based
depth-first iteration, avoiding function-call overhead.

\emph{Relationship to GCPC.}
The \texttt{GCPC} source (\cref{sec:gcpc}) adds Phase~1---GCPC
cache lookup for interior critical points---on top of the same
boundary enumeration.
CritSample uses only boundary enumeration (Phase~2),
making it faster ($\sim$39\,\textmu s vs.
$\sim$7\,ms for GCPC) while capturing 99.4\% of
the analytical endpoint-iAABB volume.

\emph{Complexity.}
\begin{equation}
  T_\text{CritSample}
  = O\!\Bigl(\prod_{j=1}^{D} |C_j|\Bigr)
  \cdot O\bigl(\text{mat-mul}\bigr),
\end{equation}
where the per-configuration cost is a single 4$\times$4 matrix
multiply (no trigonometry).

\subsection{Geometric Critical-Point Cache (GCPC)}\label{sec:gcpc}
% Code: include/sbf/envelope/gcpc.h, src/envelope/gcpc.cpp

The analytical solver's per-box cost ($\sim\!9$\,ms, IIWA14)
is dominated by combinatorial boundary enumeration.
GCPC amortises this cost across all queries for a given robot
via an eight-phase query pipeline.

\emph{One-time precomputation.}
For each active link~$\ell$, critical configurations are
enumerated within full joint limits using coordinate descent.
Each configuration is stored as a \texttt{GcpcPoint} with fields
$(\ell,\, q_\text{eff},\, A,\, B,\, C,\, R)$ and
organised into a per-link KD-tree in the joint subspace.

\emph{Symmetry reductions.}
\begin{itemize}
  \item \textbf{$q_0$ elimination.}
    The radial component $R = \sqrt{p_x^2 + p_y^2}$ and $p_z$ do
    not depend on $q_0$.
    Joint $q_0$ is not stored; at query time it is recovered via
    $q_0 = \mathrm{atan2}(-B, A)$.
  \item \textbf{$q_1$ period-$\pi$ symmetry.}
    $R(q_1) = R(q_1 + \pi)$; only $q_1 \in [0,\pi]$ is stored;
    the mirror gives the $q_1 - \pi$ candidate.
\end{itemize}

\emph{Eight-phase query pipeline.}
Given a C-space box~$B$:
\begin{enumerate}
  \item[\textbf{A.}] \textbf{Cache lookup.}
    KD-tree range query retrieves all cached points within~$B$;
    $q_0/q_1$ reconstruction establishes a strong prior bound.
  \item[\textbf{B.}] \textbf{Boundary $k\pi/2$ enumeration.}
    Boundary + interior $k\pi/2$ grid.
  \item[\textbf{C.}] \textbf{1D atan2 edge solve.}
    Per-joint sinusoidal solve via \texttt{extract\_1d\_coefficients}
    (direct DH-chain extraction, zero FK) with two-level AA pruning.
  \item[\textbf{D.}] \textbf{2D face companion-matrix solve.}
    Per-pair bilinear-trig solve via \texttt{extract\_2d\_coefficients}
    (zero FK) with AA pruning.
  \item[\textbf{D$\tfrac{1}{2}$.}] \textbf{$k\pi/2$-background cache lookup.}
    Precomputed configurations where \emph{all} background joints
    take $k\pi/2$ values are checked against the query interval
    in $O(1)$ per entry, avoiding redundant enumeration in
    Phases~E/F.
  \item[\textbf{E.}] \textbf{Pair-constrained 1D.}
    $q_i + q_j = k\pi/2$ with \texttt{extract\_2d\_coefficients}
    + constraint substitution to extract $\alpha,\beta$ sinusoidal
    coefficients (zero FK); background joints restricted to
    $\{l_j, u_j\}$ only ($k\pi/2$ values handled by Phase~D$\tfrac{1}{2}$).
  \item[\textbf{F.}] \textbf{Pair-constrained 2D.}
    $q_i + q_j = k\pi/2$ with one additional free joint.
    Precomputed DH transforms and \texttt{build\_symbolic\_poly8}
    polynomial solver; QR replaced by precomputed matrix inverse;
    incremental chain recomputation partitions background joints
    into four segments (prefix/mid12/mid23/suffix) and recomputes
    only dirty segments per combination.
  \item[\textbf{G.}] \textbf{Interior coordinate descent.}
    Multi-start, dual-mode (with and without Phase~D seeds).
\end{enumerate}
Phase~A's strong prior bound enables aggressive AA pruning in
Phases~B--G, dramatically reducing FK calls.
Phase~D$\tfrac{1}{2}$ decouples $k\pi/2$ interior values from the
boundary loop, reducing background combinations from
$\sim\!5^k$ to $2^k$ in Phases~E/F.

\emph{AA pruning (Axis-Achievability).}
For each link~$\ell$ and axis~$d$, if the current estimate
already matches the iFK outer bound, further solving on that
link/axis is skipped.
Pruning is refreshed at the start of each phase.

\emph{Precomputed $k\pi/2$-background configurations.}
During one-time cache construction, \texttt{precompute\_pair\_kpi2()}
enumerates all background-joint assignments where every background
joint takes a $k\pi/2$ value (up to 64 combinations per link pair).
The resulting critical configurations are stored in
\texttt{PairKpi2Config} records and checked at query time in
Phase~D$\tfrac{1}{2}$ via interval containment, with zero
FK cost.

\emph{Cost--tightness trade-off.}
With the eight-phase pipeline (including $k\pi/2$-background caching
and incremental chain recomputation),
GCPC produces endpoint-iAABBs at $0.89{\times}$ the cost of the
analytical solver while achieving near-identical envelope volumes
($0.121$--$0.123$\,m$^3$ for Hull16\_Grid).
The cache is \emph{scene-independent} and serialisable to disk.

\subsection{Certified BnB Verifier with Monotonicity Reduction}%
\label{sec:certified_verifier}
% Code: include/sbf/envelope/certified_verifier.h
% Code: src/envelope/certified_verifier.cpp

The IFK-clamped analytical solver~(\cref{sec:ifk_cert}) ensures
safety but relies on heuristic coordinate descent for $|F|\geq 3$,
leaving a completeness gap for high-dimensional faces.
We close this gap with a \emph{certified verifier} that combines
the Krawczyk interval-Newton operator with
\emph{monotonicity-based dimension reduction}.

\subsubsection{Krawczyk Verification for Low Dimensions}

For a single free joint~$j$, the derivative
$f(q_j) = -A_j\sin q_j + B_j\cos q_j$
has a known interval Krawczyk operator:
\begin{equation}\label{eq:krawczyk1d}
  K(X) = \hat x - Y f(\hat x) + (1 - Y f'(X))(X - \hat x),
\end{equation}
where $\hat x = \text{mid}(X)$ and $Y = 1/f'(\hat x)$.
By the Krawczyk theorem, $K(X) \subset \text{int}(X)$
implies a unique root in~$X$, while $K(X) \cap X = \emptyset$
certifies no root.
A branch-and-bound (BnB) engine bisects unresolved cells
until all are certified or a depth limit is reached.

The same principle extends to 2D using bilinear-trig
coefficients~(\cref{prop:bilinear}) with a 2D Krawczyk operator
and four boundary-edge 1D sub-problems.

\subsubsection{Monotonicity-Based Dimension Reduction}

For $n \geq 3$ free dimensions, direct BnB is impractical
(exponential cell count).
We observe that the partial derivative
\begin{equation}\label{eq:mono_grad}
  \frac{\partial p_{\ell,d}}{\partial q_j}
  = -\alpha_{j,d}\sin q_j + \beta_{j,d}\cos q_j
\end{equation}
can be enclosed in an interval using the prefix--suffix
factorisation~(\cref{eq:prefix_suffix}) and certified interval
arithmetic (outward rounding):
\begin{equation}\label{eq:mono_test}
  \Bigl[\frac{\partial p_{\ell,d}}{\partial q_j}\Bigr]
  \supseteq -[\alpha_{j,d}]\cdot[\sin X_j]
  + [\beta_{j,d}]\cdot[\cos X_j].
\end{equation}
If $0 \notin [\partial p_{\ell,d}/\partial q_j]$, then
$p_{\ell,d}$ is \emph{strictly monotone} in~$q_j$ over the
entire box, and its extremum in that dimension lies on
the boundary.

\begin{proposition}[Monotonicity probability]\label{prop:mono_prob}
For a box of half-width~$w$, the probability that dimension~$j$
is non-monotone is at most $w/\pi$ (since the unique zero modulo
$\pi$ of~\eqref{eq:mono_grad} falls in $X_j$ with probability
$\leq w/\pi$).
For $w = 0.005$\,rad and $n = 7$, the probability that all
dimensions are monotone exceeds $98.9\%$.
\end{proposition}

\begin{theorem}[Monotone corner optimality]\label{thm:mono_corner}
If $p_{\ell,d}$ is strictly monotone in every free
dimension~$j$ over box~$B$, then the global extremum is
attained at the corner determined by the gradient signs:
\begin{equation}\label{eq:corner}
  q_j^{\min} = \begin{cases}
    X_j^{\mathrm{lo}} & \text{if } \partial p/\partial q_j > 0, \\
    X_j^{\mathrm{hi}} & \text{otherwise},
  \end{cases}
\end{equation}
and symmetrically for the maximum.
\end{theorem}
\begin{proof}
Starting from any $\bm{q} \in B$, replacing each coordinate
$q_j$ with $q_j^{\min}$ decreases $p_{\ell,d}$ by strict
monotonicity.  Iterating over all dimensions yields
$p_{\ell,d}(\bm{q}) \geq p_{\ell,d}(\bm{q}^{\min})$.
\end{proof}

The verifier \texttt{verify\_nd\_mono} proceeds as:
\begin{enumerate}
  \item Compute one interval-FK pass to obtain
    $[\alpha_{j,d}], [\beta_{j,d}]$ for each free dimension.
  \item Test monotonicity via~\eqref{eq:mono_test}; partition
    free dimensions into monotone set~$M$ and non-monotone set~$\bar{M}$.
  \item If $|\bar{M}| = 0$ (fast path, $\geq 98\%$ of narrow boxes):
    evaluate two scalar FK calls at the min/max corners
    (\cref{thm:mono_corner}) to obtain exact extremal bounds.
  \item If $|\bar{M}| = 1$ or $2$: fix monotone dimensions to
    their extremal boundaries and dispatch to
    \texttt{verify\_1d} / \texttt{verify\_2d}.
  \item If $|\bar{M}| \geq 3$ (rare): apply Gauss--Seidel
    component-wise 1D Krawczyk contraction, then fall back to BnB.
\end{enumerate}

\begin{theorem}[Certified completeness]\label{thm:cert_complete}
The monotonicity test~\eqref{eq:mono_test} is sound:
it uses interval-certified coefficient enclosures
(IntervalCoeff1D with outward rounding) and
$[\sin],[\cos]$ interval extensions.
When all free dimensions are certified monotone,
\cref{thm:mono_corner} guarantees that the returned
corner values are the true global extrema of~$p_{\ell,d}$
over~$B$, achieving theoretical completeness with
zero over-approximation.
\end{theorem}

% ==============================================================================
\section{Sparse Voxel BitMask}\label{sec:voxel}
% ==============================================================================

% Code: include/sbf/voxel/bit_brick.h, include/sbf/voxel/voxel_grid.h,
%       include/sbf/voxel/hull_rasteriser.h, src/voxel/hull_rasteriser.cpp

The iAABB envelope represents each link as a single
axis-aligned box derived from its endpoint-iAABB pair.
While cheap, this representation wastes volume at joints with
large angular range.
We replace per-link iAABBs with a
\emph{Sparse Voxel BitMask} that tightly fills the true
swept-capsule geometry.

\subsection{BitBrick Layout and Spatial Hashing}\label{sec:bitbrick}

The workspace is decomposed into $8^3{=}512$-voxel tiles
(\emph{BitBricks}), each stored as $8{\times}\mathtt{uint64}$
(64~bytes, one cache line).
Tiles are addressed by an integer $(b_x,b_y,b_z)$ triple
and stored in a hash map (FNV-1a keyed).

\subsection{16-Point Convex Hull Rasterisation}\label{sec:hull16}

Given proximal endpoint box~$B_1$ and distal endpoint box~$B_2$
for a link, the swept volume is the convex hull:
\begin{equation}\label{eq:hull16}
  S_\ell = \mathrm{Conv}\bigl(
    \mathrm{corners}(B_1) \cup \mathrm{corners}(B_2)
  \bigr),
\end{equation}
with at most 16 vertices.
At parameter~$t \in [0,1]$, the cross-section
$B(t) = (1{-}t)\,B_1 + t\,B_2$ is an iAABB with bounds
\emph{linear} in~$t$:
\begin{equation}\label{eq:bt_linear}
  \mathrm{lo}_a(t) = b_{1,\mathrm{lo}}^a + t\,\Delta\mathrm{lo}^a, \qquad
  \mathrm{hi}_a(t) = b_{1,\mathrm{hi}}^a + t\,\Delta\mathrm{hi}^a.
\end{equation}

\subsection{Turbo Scanline Algorithm}\label{sec:turbo}

We eliminate per-voxel distance evaluations by exploiting the
linear structure of~\eqref{eq:bt_linear}.

\emph{YZ $t$-range solver.}
For each $(y,z)$ scanline, the feasible $t$-range
$[t_{\mathrm{lo}}, t_{\mathrm{hi}}]$ is a single interval
solved in $O(1)$.

\emph{Conservative $x$-fill.}
Since $\mathrm{lo}_x(t)$ and $\mathrm{hi}_x(t)$ are linear in~$t$,
extremes over $[t_\mathrm{lo}, t_\mathrm{hi}]$ are at the endpoints.

\emph{Batch bit fill.}
Contiguous $x$-ranges within a BitBrick are filled via a single
64-bit OR with a precomputed mask.

\emph{Brick-batch processing with tile-level early-out.}
Instead of iterating over individual $(c_y,c_z)$ cells, the YZ loop
is grouped into $8{\times}8$ brick tiles.
For each tile $(b_y,b_z)$, a single centre-point test with an
inflated radius $r_{\mathrm{test}} = r_{\mathrm{eff}} + 3.5\sqrt{2}\,\delta$
conservatively rejects tiles entirely outside the hull,
skipping up to 64 scanlines.
Within a tile, $x$-direction bitmasks accumulate into stack-local
\texttt{BitBrick} arrays (one per $x$-brick); the global hash map is
updated only once per tile, reducing hash probes from
$O(n_{\mathrm{cells}} \cdot n_{x\text{-bricks}})$ to
$O(n_{yz\text{-tiles}} \cdot n_{x\text{-bricks}})$.

\emph{Adaptive delta.}
For large bounding volumes, the total number of YZ scanlines
$\sum_\ell \lceil \mathrm{ext}_y^\ell / \delta \rceil
         \lceil \mathrm{ext}_z^\ell / \delta \rceil$
can grow without bound.
When the estimated count exceeds a threshold $N_{\max}{=}40{,}000$,
the resolution is automatically coarsened:
$\delta \leftarrow \max\!\bigl(\delta,\;
\sqrt{\textstyle\sum_\ell \mathrm{ext}_y^\ell \cdot
\mathrm{ext}_z^\ell \;/\; N_{\max}}\bigr)$.

Overall complexity: $\calO(n_y \cdot n_z)$, $O(1)$ per scanline.

\subsection{Voxel Collision and Merge}\label{sec:voxel_ops}

\emph{Collision detection.}
A single 64-bit AND per word with early exit:
\begin{equation}
  \begin{aligned}
    \text{collide}(G_R, G_O) \;\Leftrightarrow\;&\exists\, \bm{b}: \\
    &\bigl(G_R[\bm{b}].\mathrm{words}[k]\\
    &\qquad\mathbin{\&}\; G_O[\bm{b}].\mathrm{words}[k]\bigr) \ne 0.
  \end{aligned}
\end{equation}

\emph{Zero-loss merge.}
Let $P,L,R$ denote the parent, left-child, and right-child grids.
\begin{equation}
  G_P[\bm{b}] = G_L[\bm{b}] \mathbin{|} G_R[\bm{b}].
\end{equation}
Unlike iAABB union, voxel OR preserves exact shape with zero volume
inflation.

% ==============================================================================
\section{SafeBoxForest Architecture}\label{sec:sbf}
% ==============================================================================

% Code: include/sbf/envelope/endpoint_source.h, src/envelope/endpoint_source.cpp
%       include/sbf/forest/, src/forest/

\subsection{Lifelong Envelope Cache Tree (LECT)}\label{sec:lect}

The LECT is a lazy KD-tree over C-space that caches joint-origin
interval vectors and derives envelopes on demand.

\emph{Structure.}
Each node~$v$ stores a C-space box $B_v$ and $L$ joint-origin
interval vectors.
The link envelope is derived on demand as
\begin{equation}\label{eq:lect_iaabb}
  \iAABB_\ell(B_v)=\mathrm{bbox}\!\bigl([\bm{p}_{\ell-1}]_{B_v},\,[\bm{p}_\ell]_{B_v}\bigr)\oplus r_\ell.
\end{equation}

\emph{Lazy splitting.}
The LECT starts with only the root node and splits on demand.
When splitting along dimension~$d$, the prefix $[\bT^{0:d-1}]$ is
shared; only $[\bT^{d:D}]$ is recomputed (43\% FK saving on average).
A per-depth heuristic (\texttt{BEST\_TIGHTEN}) selects the dimension
whose bisection minimises the minimax link-envelope volume:
for each candidate~$d$, incremental FK is computed for both child
intervals, and
$\text{metric}(d) = \max\bigl(\sum_i V_i^L,\; \sum_i V_i^R\bigr)$
is evaluated; the dimension with the smallest metric is chosen.

\emph{IFK fast path with partial inheritance.}
% Code: src/forest/lect.cpp — compute_envelope()
When the endpoint source is IFK and a valid incremental FK state is
available (which is the common case during \texttt{split\_leaf} and
\texttt{find\_free\_box} descent), endpoint coordinates are extracted
directly from the FK prefix matrices
$\bm{p}_\ell = \bT^{0:\ell}[0{:}2,3]$,
completely bypassing \texttt{compute\_endpoint\_iaabb()} and its
redundant full FK recomputation.
Furthermore, when the split dimension~$d$ and the parent node are
known, links whose frame index satisfies $V+1 \leq d$ are
\emph{unchanged} by the split; their paired endpoint data is copied
from the parent via \texttt{memcpy}, and only the affected links are
extracted from the FK state.
This reduces per-node FK cost by 40--45\% compared to the
non-fast-path baseline.

\emph{SoA layout.}
All per-node data is stored in Structure-of-Arrays layout with
mmap lazy-load and incremental append for warm start.

\emph{Cache portability.}
The LECT depends only on robot kinematics (DH parameters and link
radii), not on obstacles.
A robot fingerprint hash validates cache on load.

\emph{Runtime source override.}
The LECT cache file records the endpoint source used during
cache generation (e.g.\ \texttt{CritSample}).
At runtime, the planner may override this with a different source
(e.g.\ \texttt{IFK}) via \texttt{set\_ep\_config()}.
Subsequent \texttt{expand\_leaf} calls use the overridden source
for newly created nodes, while previously cached nodes retain their
original data in a separate channel.
The dual-channel architecture (CH\_SAFE for IFK, CH\_UNSAFE for
sampled sources) allows mixing: collision queries prefer the tighter
channel when available, falling back to the other.
This enables loading a CritSample-generated cache for broad coverage
while using IFK's $140{\times}$ faster incremental path for
runtime expansion.

\subsection{Endpoint Source Dispatcher}\label{sec:endpoint_dispatch}
% Code: include/sbf/envelope/endpoint_source.h

The \texttt{EndpointSource} module provides a unified dispatch
(\texttt{compute\_endpoint\_iaabb()}) that selects among the
four sources based on \texttt{EndpointSource} enum:
\texttt{IFK}, \texttt{CritSample}, \texttt{Analytical}, \texttt{GCPC}.
The output is always \texttt{EndpointIAABBResult}
containing $n_\text{endpoints} \times 6$ floats.

\emph{Quality ordering (tightness).}
Sources are ranked by envelope tightness:
$\mathrm{IFK}(0) < \mathrm{CritSample}(1) <
\mathrm{Analytical}(2) < \mathrm{GCPC}(3)$.
GCPC uses the analytical solver for boundary phases and cache
lookup for interior critical points, producing strictly tighter
bounds than standalone Analytical whose Phase~3 coordinate descent
may miss interior extrema.
A cached result of quality $\geq$ the requested quality can serve
the request without recomputation.

\emph{Safety ordering.}
The safety ordering is \emph{reversed} with respect to tightness.
IFK is the only source that produces \emph{certified} envelopes
(conservative over-approximation $\supseteq$ true occupancy).
Analytical and GCPC compute near-exact envelopes by solving for
true extrema; they are tighter than IFK and safer than CritSample,
but because Phase~3 is heuristic ($|F| \geq 3$), they do not
formally guarantee $\supseteq$ true occupancy.
In practice, collisions due to this incompleteness are
extremely rare.
CritSample, using only boundary enumeration, has the
largest (though still small) completeness gap.

\subsection{Forest Growth and FFB}\label{sec:forest_growth}

Boxes are grown incrementally by placing seeds on the boundary
faces of existing boxes and calling FFB to descend the LECT.
We compare two growth strategies:

\emph{Wavefront BFS.}
Each start/goal configuration anchors a subtree.
Subtrees maintain independent BFS queues; a round-robin scheduler
pops one parent box per subtree per round, expands all boundary
face seeds via \texttt{\_generate\_boundary\_seeds()}, and calls
FFB on each.
When all BFS queues are exhausted, a guided-random fallback
samples from unoccupied LECT regions to start new subtrees.
Growth terminates at \texttt{max\_boxes} or after
\texttt{max\_consecutive\_miss} failures.

\emph{Multi-tree RRT.}
Each subtree independently extends toward random
or goal-biased targets.
The nearest box within the subtree is identified by linear scan;
\texttt{rrt\_snap\_to\_face()} selects the face most aligned with
the growth direction (maximum $\langle \mathbf{d}_k, \mathbf{n}_f \rangle$)
and places the seed as a $0.7/0.3$ blend of the directional
projection and a random point on that face.
A boundary fallback fires when all per-tree attempts are
exhausted.

\emph{FFB (Find Free Box)} descends the LECT with two-stage collision
checking (iAABB early-out, then hull-16 refinement against
pre-rasterised scene grid).

\emph{Scene pre-rasterisation.}
Obstacle iAABBs are rasterised into a shared \texttt{VoxelGrid} once.
Subsequent queries reduce to hash-map intersection with bitwise AND
($2272{\times}$ speedup over per-query grid construction).

\emph{Growth strategy comparison (2-DOF, 54 configs $\times$ 3 scenes
$=$ 162 runs).}
Wavefront BFS and RRT achieve comparable MC coverage
($\approx 0.85$ at \texttt{ffb\_min\_edge}${}=0.05$),
but Wavefront is $50$--$500{\times}$ faster
(10--50\,ms vs.\ 300--5000\,ms per forest).
\texttt{ffb\_min\_edge} is the dominant parameter:
reducing it from $0.10$ to $0.03$ raises coverage from $77\%$
to $93\%$ at the cost of $4.6{\times}$ more boxes.
Wavefront is insensitive to \texttt{guided\_sample\_ratio}
and \texttt{n\_edge\_samples} (std $< 0.01$).
The composite score (coverage, box count, connectivity, time)
favours Wavefront across all tested configurations.

\subsection{Tree-Cached Greedy Coarsening}\label{sec:coarsen}

Coarsening attempts to merge sibling boxes using zero-loss voxel OR:
$G_{v} = G_{v_L} \mathbin{|} G_{v_R}$.
The quality ratio
$\rho(v) = \mathrm{vol}(G_v) / \max(\mathrm{vol}(G_{v_L}), \mathrm{vol}(G_{v_R}))$
must satisfy $\rho \leq \rho_\text{max}$.

\subsection{GCS Graph and Incremental Update}\label{sec:gcs_graph}

Two boxes are adjacent when their C-space hyperrectangles overlap
in every dimension.
Adjacency testing uses chunked vectorised broadcasting.
Incremental scene changes invalidate only affected boxes;
the LECT cache is scene-independent.

% ==============================================================================
\section{GCS Path Planning}\label{sec:planner}
% ==============================================================================

\subsection{GCS Planning Integration}\label{sec:gcs_planning}

Each forest box becomes a GCS vertex.
Transition constraints: $\bm{q}_{ij} \in B_i \cap B_j$.
The SOCP minimises $\sum \|\bm{q}_{i,\text{in}} - \bm{q}_{i,\text{out}}\|^2$.
Dijkstra pre-solve restricts the SOCP to a topological path
plus 1-hop neighbourhood.

\subsection{Path Post-Processing}\label{sec:post_process}

Shortcutting removes unnecessary waypoints via line-of-sight tests
within box chains.
B-spline smoothing fits the remaining waypoints,
constrained to the union of traversed boxes.
No point-wise collision checking is needed.

% ==============================================================================
\section{Experiments}\label{sec:experiments}
% ==============================================================================

% Auto-generated tables and figures:
%   python python/scripts/gen_tables.py --data-dir results/paper --output-dir paper/generated
%   python python/scripts/gen_figures.py --data-dir results/paper --output-dir paper/fig

All experiments use a 7-DOF KUKA IIWA14 manipulator with
$L_\text{act}{=}4$ active links and six benchmark scenes
(three 2-DOF planar and three 7-DOF Panda arm).
Timings are measured in Release (MSVC 2022, /O2) on an Intel Core~i7.
Each experiment runs 10 seeded trials per (planner, scene) pair;
we report mean~$\pm$~std and use the Wilcoxon signed-rank test
($\alpha{=}0.05$) for pairwise significance.

\subsection{Envelope Volume and Timing}\label{sec:exp_pipeline}

We evaluate the $4{\times}3$ Cartesian product of four
\emph{endpoint-iAABB sources} and three \emph{envelope representations}
(12~pipeline configurations).

\begin{itemize}
  \item \textbf{Endpoint-iAABB Sources.}
    \emph{IFK}: interval FK (${\sim}22$\,\textmu s);
    \emph{CritSample}: boundary $k\pi/2$ enumeration with
    precomputed DH transforms (${\sim}39$\,\textmu s);
    \emph{AnalyticalCrit}: five-phase solver with DH direct
    coefficient extraction (${\sim}9$\,ms);
    \emph{GCPC}: eight-phase cache query with $k\pi/2$ cache
    and incremental chain (${\sim}7$\,ms).
  \item \textbf{Envelope Representations.}
    \emph{LinkIAABB}\footnote{\LinkIAABBDef}: radius-inflated bounding box of each link's
    endpoint-iAABB pair (or the union of its subdivided sub-link boxes);
    \emph{LinkIAABB\_Grid}: link-iAABBs rasterised into a grid;
    \emph{Hull16\_Grid}: 16-point hull rasterised directly from
    endpoint-iAABBs.
\end{itemize}

% Auto-generated by gen_paper_tables.py
\begin{table}[t]
\centering
\caption{Envelope volume across pipeline configurations
  (500 random boxes, width $w{=}0.2$\,rad).
  $V$ = mean volume;
  Ratio = $V / V_{\text{IFK-LinkIAABB}}$.}
\label{tab:envelope_volume}
\setlength{\tabcolsep}{4pt}
\small
\begin{tabular}{@{}ll cc cc c@{}}
\toprule
 & & \multicolumn{2}{c}{2-DOF} & \multicolumn{2}{c}{7-DOF (Panda)} & \\
\cmidrule(lr){3-4}\cmidrule(lr){5-6}
Source & Envelope & $V$ & Ratio & $V$ & Ratio & Cert.\ \\
\midrule
  IFK & LinkIAABB & 0.4435 & 100.0\% & 0.794 & 100.0\% & \checkmark \\
   & LinkIAABB\_Grid & 0.4330 & 97.6\% & 0.518 & 65.3\% & \checkmark \\
   & Hull16\_Grid & 0.3066 & 69.1\% & 0.435 & 54.8\% & \checkmark \\
  \addlinespace
  CritSample & LinkIAABB & 0.3845 & 86.7\% & 0.259 & 32.7\% &  \\
   & LinkIAABB\_Grid & 0.3587 & 80.9\% & 0.193 & 24.3\% &  \\
   & Hull16\_Grid & 0.2470 & 55.7\% & \textbf{0.117} & 14.8\% &  \\
  \addlinespace
  Analytical & LinkIAABB & 0.3799 & 85.7\% & 0.253 & 31.9\% &  \\
   & LinkIAABB\_Grid & 0.3701 & 83.5\% & 0.201 & 25.4\% &  \\
   & Hull16\_Grid & 0.2444 & 55.1\% & 0.121 & 15.3\% &  \\
  \addlinespace
  GCPC & LinkIAABB & 0.3847 & 86.7\% & 0.257 & 32.4\% &  \\
   & LinkIAABB\_Grid & 0.3657 & 82.4\% & 0.197 & 24.9\% &  \\
   & Hull16\_Grid & 0.2465 & 55.6\% & 0.123 & 15.5\% &  \\
\bottomrule
\end{tabular}
\end{table}

Key observations:
\begin{enumerate}
\item \textbf{CritSample closes the tightness gap at low cost.}
  CritSample+LinkIAABB is $2.6{\times}$ tighter
  than IFK+LinkIAABB at only $0.278$\,ms cold cost.

\item \textbf{Hull16\_Grid provides the tightest envelopes.}
  CritSample+Hull16\_Grid is $42\%$ tighter
  than CritSample+LinkIAABB.

\item \textbf{AnalyticalCrit $\approx$ GCPC.}
  Both produce near-identical volumes across all representations
  (within $1.4\%$), confirming GCPC cache fidelity.
  GCPC is $1.1{\times}$ faster.

\item \textbf{Warm mode eliminates source cost.}
  Cached endpoint-iAABBs ($1{,}632$\,B/node) reduce source cost
  to sub-microsecond.
\end{enumerate}

% Auto-generated by gen_paper_tables.py
\begin{table}[t]
\centering
\caption{Per-box computation time ($\mu$s) on 7-DOF Panda
  (50\,k samples).
  EP = endpoint-iAABB source; Env = envelope rasterisation;
  Speedup = Analytical-LinkIAABB / Total.}
\label{tab:timing}
\setlength{\tabcolsep}{3.5pt}
\small
\begin{tabular}{@{}ll rrr r@{}}
\toprule
Source & Envelope & EP\,($\mu$s) & Env\,($\mu$s) & Total\,($\mu$s) & Speedup \\
\midrule
  IFK & LinkIAABB & 4.4 & 0.01 & \textbf{4.4} & 3,528$\times$ \\
   & LinkIAABB\_Grid & 4.3 & 8.6 & 13 & 1,198$\times$ \\
   & Hull16\_Grid & 4.3 & 22 & 27 & 579$\times$ \\
  \addlinespace
  CritSample & LinkIAABB & 10 & 0.00 & 10 & 1,478$\times$ \\
   & LinkIAABB\_Grid & 11 & 4.1 & 15 & 1,047$\times$ \\
   & Hull16\_Grid & 11 & 13 & 24 & 646$\times$ \\
  \addlinespace
  Analytical & LinkIAABB & 15,432 & 0.4 & 15,432 & 1.0$\times$ \\
   & LinkIAABB\_Grid & 15,052 & 6.3 & 15,058 & 1.0$\times$ \\
   & Hull16\_Grid & 14,886 & 21 & 14,907 & 1.0$\times$ \\
  \addlinespace
  GCPC & LinkIAABB & 13,013 & 0.2 & 13,013 & 1.2$\times$ \\
   & LinkIAABB\_Grid & 13,036 & 6.7 & 13,043 & 1.2$\times$ \\
   & Hull16\_Grid & 13,023 & 22 & 13,045 & 1.2$\times$ \\
\bottomrule
\end{tabular}
\end{table}

\emph{Recommended configuration.}
\textbf{CritSample + Hull16\_Grid} for online planning
(tightest, $0.278$\,ms cold).
\textbf{GCPC + Hull16\_Grid} for accuracy validation
($1.1{\times}$ faster than Analytical with quality level~$3$).
\textbf{IFK + Hull16\_Grid} when formal safety certification
is required (conservative but provably collision-free).

\subsection{End-to-End Planning}\label{sec:exp_e2e}

We evaluate the 12~SBF pipeline configurations across all
benchmark scenes in a full build--plan--smooth cycle.

% Auto-generated by gen_paper_tables.py
\begin{table*}[t]
\centering
\caption{End-to-end planning: success rate (\%) / mean time (ms)
  across 12 SBF configurations and benchmark scenes
  (30 seeded trials each).}
\label{tab:e2e}
\setlength{\tabcolsep}{3pt}
\small
\begin{tabular}{@{}llcccccc@{}}
\toprule
Source & Envelope & 2D-S & 2D-N & 2D-C & P-Tab & P-Shf & P-Multi \\
\midrule
  IFK & LinkIAABB & 100/11.2 & 100/38.4 & 100/11.2 & 100/9,095 & 0/--- & 100/8,274 \\
   & LinkIAABB\_Grid & 100/11.2 & 100/39.0 & 100/11.3 & 100/9,061 & 0/--- & 100/8,331 \\
   & Hull16\_Grid & 100/11.2 & 100/39.0 & 100/11.2 & 100/9,181 & 0/--- & 100/8,375 \\
  \addlinespace
  CritSample & LinkIAABB & 100/11.0 & 100/42.5 & 100/11.9 & 100/13.9s & 0/--- & 100/14.9s \\
   & LinkIAABB\_Grid & 100/11.1 & 100/42.8 & 100/11.6 & 100/13.0s & 0/--- & 100/15.0s \\
   & Hull16\_Grid & 100/11.1 & 100/43.0 & 100/11.2 & 100/12.0s & 0/--- & 100/14.4s \\
  \addlinespace
  \addlinespace
\bottomrule
\end{tabular}
\end{table*}

\begin{figure}[t]
  \centering
  \IfFileExists{fig/fig2_e2e_heatmap.pdf}{%
    \includegraphics[width=\columnwidth]{fig/fig2_e2e_heatmap.pdf}%
  }{\fbox{\parbox{\columnwidth}{\centering\vspace{2cm}Fig.~2 pending\vspace{2cm}}}}
  \caption{End-to-end planning time (colour) and success rate
  (annotation) across pipeline configurations and scenes.}
  \label{fig:e2e_heatmap}
\end{figure}

\subsection{Baseline Comparison}\label{sec:exp_baselines}

We compare the recommended SBF configuration
(CritSample + Hull16\_Grid + Dijkstra) against
six sampling-based baselines:
RRT, RRT-Connect, RRT$^*$, Informed-RRT$^*$, BiRRT$^*$, and PRM.
Additionally, we compare against IRIS-NP-GCS and C-IRIS-GCS
(where Drake is available).

% Auto-generated by gen_paper_tables.py
\begin{table}[t]
\centering
\caption{Baseline comparison across benchmark scenes (30 trials, mean~$\pm$~std).}
\label{tab:baselines}
\small
\begin{tabular}{@{}lccc@{}}
\toprule
Planner & Success\,(\%) & Time\,(ms) & Path Length \\
\midrule
  SBF-IFK-LinkIAABB & 75 & 2135.3$\pm$2988.7 & 1.99 \\
\bottomrule
\end{tabular}
\end{table}

\begin{figure}[t]
  \centering
  \IfFileExists{fig/fig3_baseline_pareto.pdf}{%
    \includegraphics[width=\columnwidth]{fig/fig3_baseline_pareto.pdf}%
  }{\fbox{\parbox{\columnwidth}{\centering\vspace{2cm}Fig.~3 pending\vspace{2cm}}}}
  \caption{Pareto front: planning time vs.\ path length.
  Error bars = $\pm 1$ std over 10 trials.
  SBF achieves the fastest planning time while maintaining
  competitive path quality.}
  \label{fig:baseline_pareto}
\end{figure}

SBF achieves the highest success rate and fastest mean planning time
across all scenes, with path lengths comparable to RRT$^*$.
The key advantage is deterministic coverage: once the forest is built,
queries are guaranteed to find a path if one exists within the
box connectivity graph.

\subsection{LECT Cache Effectiveness}\label{sec:exp_lect}

We measure the per-node expansion cost inside the LECT tree
to quantify the benefit of Z4 caching and partial FK inheritance.
\emph{Cold} denotes root node construction (no parent or cache);
\emph{Hot} denotes deeper expansions (depth${\geq}2$) that
exploit the Z4 endpoint-iAABB cache and incremental FK.
Three random intervals are generated per width setting;
timings are averaged over all deeper nodes.

% Auto-generated from bench_endpoint_iaabb — LECT expansion cold/hot timing
% Cold = root node construction (no cache/parent).
% Hot  = deeper expansion (Z4 cache, partial inheritance, incremental FK).
% 3 random intervals per width, averaged. Release MSVC 2022 / Intel Core i7.
\begin{table*}[t]
\centering
\caption{LECT expansion timing: cold start vs.\ hot (cached) expansion
across interval widths and endpoint-iAABB sources.
Cold = root construction (no parent/cache);
Hot = depth${\geq}2$ expansion with Z4 cache and partial FK inheritance.
Three random intervals per width, averaged.}
\label{tab:lect_expansion}
\small
\begin{tabular}{@{}ll rr r rr rr@{}}
\toprule
& & \multicolumn{3}{c}{Timing} & \multicolumn{2}{c}{Volume (m$^3$)} & \multicolumn{2}{c}{Extent (m)} \\
\cmidrule(lr){3-5}\cmidrule(lr){6-7}\cmidrule(lr){8-9}
Source & Width\,(rad) & Cold\,($\mu$s) & Hot\,($\mu$s) & Speedup
       & Root & Leaf & Root & Leaf \\
\midrule
\multicolumn{9}{l}{\textit{Panda 7-DOF ($n_\text{joints}{=}7$, $n_\text{active}{=}5$, tool frame)}} \\
\midrule
\multirow{4}{*}{IFK}
  & 0.10 &    73.8 &   29.6 & 2.49$\times$ & 0.016 & 0.002 & 0.084 & 0.043 \\
  & 0.30 &    57.9 &   29.2 & 1.98$\times$ & 0.501 & 0.071 & 0.263 & 0.132 \\
  & 0.50 &    70.7 &   30.6 & 2.31$\times$ & 2.707 & 0.389 & 0.459 & 0.234 \\
  & 1.00 &    61.6 &   31.3 & 1.97$\times$ & 26.30 & 3.726 & 0.967 & 0.491 \\
\midrule
\multirow{4}{*}{CritSample}
  & 0.10 &   296.7 &  269.1 & 1.10$\times$ & 0.002 & $<$0.001 & 0.047 & 0.025 \\
  & 0.30 &   291.1 &  271.3 & 1.07$\times$ & 0.062 & 0.011 & 0.145 & 0.078 \\
  & 0.50 &   328.4 &  289.7 & 1.13$\times$ & 0.290 & 0.051 & 0.246 & 0.132 \\
  & 1.00 &   387.7 &  306.3 & 1.27$\times$ & 2.096 & 0.421 & 0.476 & 0.267 \\
\midrule
\multirow{4}{*}{Analytical}
  & 0.10 & 14\,791 & 15\,959 & 0.93$\times$ & 0.002 & $<$0.001 & 0.047 & 0.025 \\
  & 0.30 & 15\,687 & 15\,855 & 0.99$\times$ & 0.062 & 0.011 & 0.145 & 0.078 \\
  & 0.50 & 15\,217 & 15\,049 & 1.01$\times$ & 0.294 & 0.051 & 0.246 & 0.132 \\
  & 1.00 & 14\,930 & 15\,033 & 0.99$\times$ & 2.198 & 0.430 & 0.482 & 0.269 \\
\midrule
\multirow{4}{*}{GCPC}
  & 0.10 & 12\,606 & 13\,283 & 0.95$\times$ & 0.002 & $<$0.001 & 0.047 & 0.025 \\
  & 0.30 & 13\,564 & 13\,314 & 1.02$\times$ & 0.062 & 0.011 & 0.145 & 0.078 \\
  & 0.50 & 12\,995 & 12\,444 & 1.04$\times$ & 0.294 & 0.051 & 0.246 & 0.132 \\
  & 1.00 & 12\,939 & 12\,635 & 1.02$\times$ & 2.198 & 0.430 & 0.482 & 0.269 \\
\midrule
\multicolumn{9}{l}{\textit{2-DOF planar ($n_\text{joints}{=}2$, $n_\text{active}{=}1$)}} \\
\midrule
\multirow{5}{*}{IFK}
  & 0.10 &    39.9 &    5.8 & 6.83$\times$ & --- & --- & 0.020 & $<$0.001 \\
  & 0.30 &    19.3 &    6.3 & 3.06$\times$ & --- & --- & 0.059 & 0.001 \\
  & 0.50 &    27.6 &    6.3 & 4.35$\times$ & --- & --- & 0.096 & 0.002 \\
  & 1.00 &    35.3 &    6.6 & 5.33$\times$ & --- & --- & 0.190 & 0.003 \\
  & 2.00 &    61.8 &    6.8 & 9.11$\times$ & --- & --- & 0.390 & 0.007 \\
\midrule
\multirow{5}{*}{GCPC}
  & 0.10 &    34.7 &   12.3 & 2.81$\times$ & --- & --- & 0.020 & $<$0.001 \\
  & 0.30 &    37.4 &   11.7 & 3.19$\times$ & --- & --- & 0.059 & 0.001 \\
  & 0.50 &    39.8 &   13.2 & 3.01$\times$ & --- & --- & 0.096 & 0.002 \\
  & 1.00 &    43.0 &   15.3 & 2.82$\times$ & --- & --- & 0.190 & 0.003 \\
  & 2.00 &    50.8 &   14.5 & 3.50$\times$ & --- & --- & 0.390 & 0.007 \\
\bottomrule
\end{tabular}
\end{table*}


Key observations:
\begin{enumerate}
  \item \textbf{IFK benefits most from caching.}
    On the 7-DOF Panda, IFK achieves $2.0$--$2.5{\times}$ hot speedup
    because the cached Z4 endpoints eliminate repeated interval FK
    evaluations.
    On the 2-DOF planar robot, the speedup reaches $3$--$9{\times}$
    as the per-node overhead shrinks to ${\sim}6$\,\textmu s.

  \item \textbf{CritSample shows moderate benefit.}
    Hot CritSample is $1.1$--$1.3{\times}$ faster,
    limited by the per-node boundary enumeration cost
    (${\sim}270$\,\textmu s) that caching cannot eliminate.

  \item \textbf{Analytical and GCPC: computation dominates.}
    Both sources spend $12$--$16$\,ms per node regardless of
    cache state (speedup ${\approx}1.0{\times}$).
    Their critical-point solvers recompute from scratch at each node;
    Z4 caching saves only the final endpoint-iAABB storage, not
    the solver cost.

  \item \textbf{IFK envelopes are looser.}
    IFK root volumes are $7$--$12{\times}$ larger than
    CritSample/Analytical/GCPC at the same width,
    confirming the volume--speed trade-off:
    IFK is fast and certified but conservative;
    CritSample provides near-analytical tightness at sub-millisecond cost.

  \item \textbf{Narrower intervals benefit more from subdivision.}
    Leaf-to-root volume ratios decrease with width:
    at $w{=}0.10$\,rad, leaf volume is ${\sim}8{\times}$ smaller
    than root volume, justifying deeper LECT expansion for tight boxes.
\end{enumerate}

\subsection{Scalability}\label{sec:exp_scalability}

\begin{figure*}[t]
  \centering
  \IfFileExists{fig/fig4_scalability.pdf}{%
    \includegraphics[width=\textwidth]{fig/fig4_scalability.pdf}%
  }{\fbox{\parbox{\textwidth}{\centering\vspace{2cm}Fig.~4 pending\vspace{2cm}}}}
  \caption{Scalability analysis. (a)~Planning time vs.\ DOF.
  (b)~Success rate vs.\ number of obstacles.
  (c)~Success rate vs.\ box budget.}
  \label{fig:scalability}
\end{figure*}

\emph{Envelope tightness vs.\ width.}
Envelope volume grows as $\calO(w^2)$ for narrow boxes
and saturates for wide ones.

\emph{Link subdivision tightness.}
We evaluate the union-volume metric
(Section~\ref{sec:subdivision}) for link subdivision
with $N_{\text{sub}} \in \{1,2,4,8\}$ sub-segments
across six interval widths (\cref{tab:subdivision}).
The tight sub-AABBs are inflated by per-link radius
only at grid rasterisation time, avoiding redundant
per-sub-AABB padding.
Union volume decreases monotonically with~$N_{\text{sub}}$
at all widths; the improvement is most pronounced for
narrow intervals ($37\%$ reduction from $N_{\text{sub}}{=}1$
to~$8$ at $w{=}0.10$\,rad), confirming that link subdivision
is most effective where the wrapping effect is smallest.
At wide intervals ($w{\geq}1.0$\,rad) the link envelopes
are dominated by the C-space sweep and subdivision provides
diminishing returns.

% Auto-generated by gen_paper_tables.py
\begin{table}[t]
\centering
\caption{Link-iAABB subdivision: union volume (m$^3$) decreases
  monotonically with subdivision count $N_{\text{sub}}$.
  IIWA14, CritSample source, 3 random intervals averaged.}
\label{tab:subdivision}
\small
\begin{tabular}{@{}rrrrrrr@{}}
\toprule
$N_{\text{sub}}$ & $w{=}0.10$ & $w{=}0.20$ & $w{=}0.40$ & $w{=}0.60$ & $w{=}1.00$ & $w{=}1.60$ \\
\midrule
  1 & 0.1090 & 0.1482 & 0.2740 & 0.5042 & 1.4079 & 3.2455 \\
  2 & 0.0831 & 0.1217 & 0.2469 & 0.4772 & 1.3921 & 3.2455 \\
  4 & 0.0723 & 0.1108 & 0.2352 & 0.4667 & 1.3849 & 3.2455 \\
  8 & 0.0670 & 0.1049 & 0.2301 & 0.4618 & 1.3818 & 3.2455 \\
\bottomrule
\end{tabular}
\end{table}

\emph{Voxel collision.}
Pre-rasterised scene reduces hull collision from $0.47$\,ms
to $3.3$\,ns per query ($2272{\times}$).

\emph{SBF construction.}
200 boxes covering a 10-obstacle tabletop scene built in
$<\!1$\,s cold; incremental update after perturbation
takes $<\!200$\,ms.

\emph{Growth strategy.}
Wavefront BFS vs.\ multi-tree RRT across
$54$ parameter configurations $\times$ $3$ scenes
(2-DOF, see Sec.~\ref{sec:forest_growth}):
Wavefront achieves mean coverage $0.85$ at \texttt{ffb\_min\_edge}${}=0.05$
in $32$\,ms (mean), while RRT reaches $0.85$ in $2{,}100$\,ms---a
$50$--$500{\times}$ speed gap.

% ==============================================================================
\section{Conclusion}\label{sec:conclusion}
% ==============================================================================

We presented SafeBoxForest v5, a system for rapidly covering
$\Cfree$ with certified convex boxes for GCS-based motion planning.
Key contributions:
(i)~a modular two-layer pipeline where four interchangeable
endpoint-iAABB sources feed into three envelope representations
(LinkIAABB\footnote{\LinkIAABBDef}, LinkIAABB\_Grid, Hull16\_Grid)
via a unified intermediate format;
(ii)~a five-phase analytical solver with DH-chain direct
coefficient extraction (zero-FK optimisation, $37\%$ speedup)
and provable face-$k$ completeness for $k \leq 2$;
(iii)~an eight-phase GCPC with $q_0$ elimination, $q_1$ period-$\pi$
symmetry, $k\pi/2$-background caching, incremental chain
recomputation, and AA pruning, achieving analytical-quality bounds at
$1.1{\times}$ lower cost;
(iv)~a GCPC-based CritSample combining cache lookup with boundary
$k\pi/2$ enumeration for sub-millisecond critical-point sampling;
(v)~the Sparse Voxel BitMask with turbo scanline rasterisation,
$42\%$ tighter than iAABBs with zero-loss merging;
(vi)~the LECT with scene pre-rasterisation ($2272{\times}$ collision
speedup) and hull-based coarsening; and
(vii)~a comprehensive $4{\times}3$ benchmark
establishing
CritSample+Hull16\_Grid as optimal
($0.117$\,m$^3$ volume, $0.278$\,ms cold).

Future work: end-to-end GCS planning evaluation against IRIS-NP,
AVX2/512 acceleration, higher-DOF extensions,
adaptive \texttt{ffb\_min\_edge} scheduling
(coarse early, fine late), inter-subtree connection detection,
and KD-tree acceleration for RRT nearest-box queries.

% ==============================================================================
\bibliographystyle{IEEEtran}
\bibliography{references}

\end{document}
